%---------------------------------------------------------------------
\subsection{Practitioner Consensus and Contested Topics}
\label{sec:practitioner}

\begin{folklorebox}{Practitioner Consensus: Core Hangboard Principles · \folkloreverified · ESS 1}
Consistent across major coaching resources (\textit{Training for Climbing}, \textit{Rock Climber's Training Manual}, Lattice Training, Dave MacLeod, Crimpd, \texttt{r/climbharder} threads, 2020--2025):
\begin{enumerate}[noitemsep, topsep=2pt]
  \item \textbf{No full-bodyweight hanging for beginners below ${\sim}$1--2 years or ${\sim}$V4/5.11} (traditional majority view): technique development is more efficient and tendons are not yet conditioned for isolated peak loads. \textit{Contested: a significant and growing camp — including Crimpd, Emil Abrahamsson (creator of Abrahangs, targeting climbers at ${\sim}5.11$), Lattice Training's feet-on-floor progressions, and Tyler Nelson's submaximal tissue-tolerance protocols — explicitly endorses early-stage low-intensity loading (feet on floor, $\leq40\%$ MFS) to build baseline tendon capacity before progressing to full bodyweight. The shift gained momentum on \texttt{r/climbharder} post-2022 and is reflected in current app-based training plans. The key distinction is ``no max-intensity loading'' vs.\ ``no finger loading whatsoever'': the former is near-universal consensus, the latter is increasingly contested.}
  \item \textbf{Minimum 15-min warm-up}: light cardio, shoulder/wrist mobility, and progressive large-hold hangs before any session.
  \item \textbf{Session freshness required}: hangboard before climbing or on a separate day; never after exhausting bouldering.
  \item \textbf{2--3 sessions/week with $\geq48$~h between hard sessions}. \textit{Note: this directly conflicts with Abrahamsson's twice-daily low-intensity protocol (Abrahangs), discussed under Contested Topics.}
  \item \textbf{Half-crimp or open-hand as default}: full crimp introduced only gradually and specifically.
  \item \textbf{Progress by load before reducing edge depth}: overload an 18--20~mm edge before moving to smaller holds.
  \item \textbf{Deload every 4--8~weeks}: reduce volume or intensity ${\sim}40$--50\% for one week to allow structural adaptation.
\end{enumerate}
\end{folklorebox}

\begin{bioplausbox}{Early Max-Intensity Loading Risk in Novice Climbers · \bioplausible · ESS 2}
Point 1 is clinically plausible---finger tendons adapt more slowly than muscles, and isolated peak loading before adequate conditioning likely elevates A2 rupture risk---and receives partial empirical support from Sjöman et al.\ \cite{SGJ23}, who found RFT positively associated with injury in less-experienced climbers reaching 7a$+$ ($p=0.03$). No controlled study has measured injury rates in beginner hangboard users vs.\ beginner-only climbers; the consensus threshold of 1--2~years / V4 is expert opinion, not an empirically derived value.
\end{bioplausbox}

\subsubsection*{Contested Topics}

\textbf{20~mm as the standard edge.} The community treats a 20~mm half-crimp as the default test and training edge on the grounds that it isolates the finger flexors without excessive skin-friction contribution yet still accommodates progressive loading. No controlled study directly compares edge depths as training stimuli; the recommendation is expert consensus, not an empirically derived threshold.

\textbf{Minimal-edge vs.\ added-weight training.} Community preference for minimal-edge work (specificity argument: hard climbing occurs on small holds) directly conflicts with Mundry et al.\ \cite{MSA+21}, the only RCT comparing the two methods, which found added weight on a standard edge significantly improved grip strength while edge reduction did not.

\textbf{Daily low-intensity hangs.} Abrahamsson's twice-daily submaximal protocol conflicts with the 48~h recovery rule but is mechanobiologically plausible and observationally supported \cite{GKAO24}. No RCT has tested this modality directly.

\textbf{No-hang devices.} No-hang tools (Tension Block, Grippul) offer precise weight-based loading with full finger-flexor isolation and no shoulder involvement. No controlled study compares them to traditional hangboarding; the clinical significance of absent shoulder engagement is unresolved.

\textbf{Full-crimp training.} Traditional guidance avoids full-crimp hangboarding given the high pulley loads \cite{Sch01}. Contemporary practitioners increasingly argue for progressive full-crimp exposure when it is used extensively on the wall, reasoning that tissue-specific inoculation is necessary. No RCT has tested this.
